% Part: propositional-logic
% Chapter: syntax-and-semantics
% Section: semantic-notions

\documentclass[../../../include/open-logic-section]{subfiles}

\begin{document}

\olfileid{pl}{syn}{sem}

\olsection{Gagasan Semantis}

Kita nemtokake gagasan-gagasan semantis iki:

\begin{defn} 
\begin{enumerate}
\item !!^a{formula}~$!A$ \emph{bisa dicukupi} yen kanggo sawijining
  $\pAssign{v}$, $\pSat{v}{!A}$; formula iku \emph{ora bisa dicukupi}
  yen ora ana $\pAssign{v}$ kang ndadekake $\pSat{v}{!A}$;
\item !!^a{formula}~$!A$ iku \emph{tautologi} yen $\pSat{v}{!A}$ kanggo
  kabeh !!{valuation}s~$\pAssign{v}$;
\item !!^a{formula}~$!A$ iku \emph{kontingen} yen bisa dicukupi nanging
  dudu tautologi;
\item Yen $\Gamma$ iku himpunan !!{formula}s, $\Gamma \Entails !A$
  (``$\Gamma$ ngakibatake $!A$ kanthi semantis'') yen lan mung yen
  $\pSat{v}{!A}$ kanggo saben !!{valuation}~$\pAssign{v}$ kang
  ndadekake $\pSat{v}{\Gamma}$.
\item Yen $\Gamma$ iku himpunan !!{formula}s, $\Gamma$
  \emph{bisa dicukupi} yen ana !!a{valuation}~$\pAssign{v}$ kang
  ndadekake $\pSat{v}{\Gamma}$, lan yen ora ana, $\Gamma$
  \emph{ora bisa dicukupi}.
\end{enumerate} 
\end{defn}

\begin{prob} 
 Kanggo saben papat !!{formula}s ing ngisor iki, temtokna apa formula
 iku (a)~bisa dicukupi, (b)~tautologi, lan (c)~kontingen.
\begin{enumerate}
  \item \( ( \Obj p_0 \lif ( \lnot \Obj p_1 \lif \lnot \Obj p_0 ) ) \).
  \item \( ( ( \Obj p_0 \land \lnot \Obj p_1 ) \lif ( \lnot \Obj p_0 \land \Obj p_2 )) \liff ( ( \Obj p_2 \lif \Obj p_0 ) \lif ( \Obj p_0 \lif \Obj p_1 )) \).
  \item \( ( \Obj p_0 \liff \Obj p_1 ) \lif ( \Obj p_2 \liff \lnot \Obj p_1 ) \).
  \item \( (( \Obj p_0 \liff ( \lnot \Obj p_1 \land \Obj p_2 )) \lor ( \Obj p_2 \lif ( \Obj p_0 \liff \Obj p_1 ))) \).
\end{enumerate}
\end{prob}

\begin{prop}
\ollabel{prop:semanticalfacts} 
\begin{enumerate} 
\item $!A$ iku tautologi yen lan mung yen
  $\emptyset \Entails !A$; 
\item Yen $\Gamma \Entails !A$ lan $\Gamma \Entails !A \lif !B$, mula
  $\Gamma \Entails !B$;
\item Yen $\Gamma$ bisa dicukupi, saben subhimpunan winates saka
  $\Gamma$ uga bisa dicukupi; 
\item \ollabel{def:monotonicity} Monotonisitas: yen $\Gamma \subseteq \Delta$
  lan $\Gamma \Entails !A$, mula uga $\Delta \Entails !A$;
\item \ollabel{def:Cut} Transitivitas: yen $\Gamma \Entails !A$ lan
  $\Delta \cup \{ !A\} \Entails !B$, mula $\Gamma \cup \Delta \Entails
  !B$.
\end{enumerate}
\end{prop}

\begin{proof}
Latihan.
\end{proof}

\begin{prob}
Buktekna \olref[pl][syn][sem]{prop:semanticalfacts}
\end{prob}

\begin{prop}\ollabel{prop:entails-unsat}
  $\Gamma \Entails !A$ yen lan mung yen $\Gamma \cup \{\lnot !A\}$
  ora bisa dicukupi.
\end{prop}

\begin{proof}
Latihan.
\end{proof}

\begin{prob}
Buktekna \olref[pl][syn][sem]{prop:entails-unsat}
\end{prob}

\begin{thm}[Teorema Deduksi Semantis]
  \ollabel{thm:sem-deduction} $\Gamma \Entails !A \lif !B$ yen lan mung
  yen $\Gamma \cup \{!A\} \Entails !B$.
\end{thm}

\begin{proof}
Latihan.
\end{proof}

\begin{prob}
Buktekna \olref[pl][syn][sem]{thm:sem-deduction}
\end{prob}

\end{document}
